\begin{center}
	
	\begin{figure}[h]
		\begin{center}
		\includegraphics[width=0.1\textwidth]{logos/logoUSB.png}
		\end{center}
	\end{figure}

	\textsc{UNIVERSIDAD SIMÓN BOLÍVAR}\\[0.1pt]
	\textsc{DECANATO DE ESTUDIOS PROFESIONALES}\\
	\textsc{COORDINACIÓN DE INGENIERÍA DE LA COMPUTACIÓN}

	\textsc{\textbf{EXTENSIÓN DE UN SISTEMA DE GESTIÓN DE APRENDIZAJE CON UN MÓDULO DE ENCUENTROS PRESENCIALES}}

	Realizado por: José Ricardo Pascarella Quijada\\
	Con la asesoría de:	Prof. Federico Flaviani e Ing. Christian Ament

	\textsc{\textbf{RESUMEN}}

\end{center}

Este documento presenta detalladamente el proceso de desarrollo de una extensión para un \gls{SGA} que permite el soporte de encuentros presenciales o seminarios entre estudiantes e instructores. 

El módulo apoya a tres tipos de usuarios en la gestión de los encuentros. Los administradores, que crean los cursos y asignan los instructores. Los instructores que pueden administrar la asistencia y calificación; por último, los estudiantes que seleccionan cursos a los que asistir.

Para la realización de este módulo se usó la metodología de desarrollo de \emph{software} ágil e iterativo por fases \emph{scrum}. Para el desarrollo se utilizó el conjunto de soluciones informáticas \gls{WAMP}, que consiste en Windows como plataforma de sistema operativo, Apache para el servicio web, \gls{SQL} Server con gestor de bases de datos y el lenguaje multipropósito \gls{PHP}.

Se logró mediante este proyecto que el SGA base de la compañía contenga la nueva funcionalidad, que permite ser extendida y personalizada dependiendo de las necesidades de distintos clientes. Además, el módulo se implantó con éxito en la arquitectura de un cliente que previamente poseía el servicio de \gls{SGA}.

\textbf{Palabras clave}: \gls{SGA}, \gls{PHP}, desarrollo, seminario.




